\begin{abstract}
Constraint Satisfaction Problems (CSPs) underpin numerous real-world applications including task planning, resource allocation, and scheduling. Large Language Models (LLMs) exhibit sharp accuracy degradation as CSP complexity grows; once the underlying solver encounters dozens of conflicts, performance approaches random and computational scaling techniques yield diminishing returns. We attribute this to two architectural limitations of existing neuro-symbolic and memory-augmented agent approaches: (1) repair mechanisms receive only coarse unsatisfiability signals and lack semantically grounded feedback, and (2) experience is either ignored across instances or stored as unverified natural-language prose without any mechanical quality control.

We propose \textbf{PRISM} (Paradigm-guided Reasoning with Iterative Solver Memory), a neuro-symbolic framework built around a single observation: in CSP solving, every reasoning step admits mechanical certification by an SMT solver, and this enables an end-to-end SMT-grounded feedback loop at two granularities. (i) At the \emph{paradigm level} (offline), candidate solving patterns distilled from successful trajectories are screened by Z3 along three axes—soundness via sampling, effect via non-contradiction and non-vacuousness checks ($\phi_{\text{pre}} \wedge \neg\,\text{op}$ must be SAT), and trigger precision via structured kind-set selectivity—before admission to the paradigm library. (ii) At the \emph{repair-trajectory level} (online), a typed memory records canonicalised UNSAT cores and structured repair signatures; UNSAT-core stagnation and structured loop detection trigger a four-level adaptive escalation protocol (switch target $\to$ switch type $\to$ revert from a bounded checkpoint stack $\to$ full retranslate) with a worst-case bounded LLM-call budget per puzzle.

To our knowledge PRISM is the first framework to apply SMT-based mechanical screening to LLM-induced reasoning rules at the paradigm level and to combine it with SMT-grounded escalation at the repair level. We evaluate PRISM on the ZebraLogic benchmark across six complexity scales (3$\times$5 to 6$\times$6) against seven neuro-symbolic and memory-based baselines, and we formulate a \emph{paradigm transfer rate} protocol to assess cross-scale and cross-problem generalisation of accumulated experience.
\end{abstract}
